\chapter{Language and Coordination Support for Collective Systems}
\label{chap:language-and-coordination}

% NOTE: this chapter gathers two publications. \S6.1--6.4 develop CaMiL
% ("Capabilities to Catch 'em All", under review at TOPLAS); \S6.5 develops
% ScalaTropy (COORDINATION 2026), whose background is
% \cref{sec:advanced-linguistic-abstractions}.
% The operational semantics of both papers is deliberately not reproduced:
% \cref{subsec:camil-formalization} states the soundness result and cites the
% paper for the rules, following the precedent of \cref{chap:pulverization}.

The pulverisation model of \cref{chap:pulverization} defines deployable components.
It also establishes that moving a component from one device to another leaves the program's behaviour unchanged.
It takes the components as given.
The three paradigms of \cref{chap:programming-models} produce these components in different ways.
Aggregate computing derives them from a field computation over an ensemble.
Choreographic programming derives them from a global protocol projected onto its participants.
Multitier programming derives them from placement annotations on the lexical scopes of a single program.
%
The literature developed each paradigm separately, with its own guarantees~\cite{DBLP:journals/csur/WeisenburgerWS20,DBLP:journals/csur/Casadei23}.

Few systems can be expressed within a single paradigm.
%
One application may need a structured protocol between fixed roles,
control over where each computation runs,
and coordination across an ensemble whose membership keeps changing.
%
A developer either picks one paradigm and encodes the others around it,
or combines separate frameworks and loses guarantees about their interaction.
%
Correct placement says nothing about protocol compliance,
and a self-stabilising collective computation gives no guarantee about communication across tiers.

Two language designs address this problem.
%
\camil turns each paradigm into a \emph{capability}: a value that code must hold to use the operations of that paradigm.
%
Placement types record where a value lives, and paradigms pass values to each other through them.
%
Scala~3's contextual abstractions, capture checking, and separation checking let the compiler enforce these rules,
and a soundness proof establishes the resulting discipline.
%
\camil is evaluated on 46 use cases from the literature and on three applications that mix paradigms in one program.
%
The second design, ScalaTropy, joins choreographic and multitier programming through monadic communication primitives rather than capabilities.
%
Choreographic and multitier languages generally provide point-to-point exchange and little else;
ScalaTropy gives each communication pattern its own operation.

\section{One System, Three Paradigms}
\label{sec:multiparty-motivation}

\cityevt~\cite{Farabegoli:2026:CaMiL}, the running example of this chapter, supports a public city event.
%
Participants join through a smartphone app.
During the event, they receive recommendations about nearby points of interest.
%
The system also monitors how crowded the event is and reports critical situations to the organisers.
%
Participants' positions are never collected:
the crowding estimate comes from a decentralised algorithm that reasons only about proximity.

\begin{figure}[t]
    \centering
    \def\svgwidth{0.50\textwidth}
    \input{figures/case-study-diagram.pdf_tex}
    \caption[The CityEvt architecture]{The \cityevt architecture: smartphones register participants, sense their neighbourhood, and display recommendations; edge servers compute ad hoc recommendations; the cloud stores participant data and collects analytics. Adapted from~\cite{Farabegoli:2026:CaMiL}.}
    \label{fig:cityevt-architecture}
\end{figure}

\Cref{fig:cityevt-architecture} shows the three types of hosts.
%
A cloud service stores participant data, manages registrations, and collects analytics;
smartphones register, report crowding, and display recommendations;
edge servers produce the recommendations.
%
The system has three parts,
each suited to a different paradigm: \emph{choreography}, \emph{multitier}, or \emph{collective}.

\paragraph{Participant registration.}
Registration involves two parties and a fixed sequence of steps:
the participant fills a form on the phone, the cloud assigns an identifier, and the phone reads it back.
%
Multitier programming (\cref{subsec:multitier-programming}) is designed for such fixed two-party workflows.
%
Its architecture specification names the \emph{peers} on which computation can be placed.
It also defines the \emph{ties} between peers and the cardinality of each tie.
%
In ScalaLoci~\cite{DBLP:journals/pacmpl/WeisenburgerKS18}, the \cityevt architecture reads as follows.

\begin{lstlisting}[
  style=scala3,
  escapeinside={(*@}{@*)},
  caption={The CityEvt architecture as a ScalaLoci peer specification.},
  label={lst:scalaloci-architecture},
]
@peer type Cloud <: { type Tie <: Multiple[Phone] & Single[Edge] }(*@\label{line:cloud-peer}@*)
@peer type Edge  <: { type Tie <: Multiple[Phone] & Single[Cloud] }(*@\label{line:edge-peer}@*)
@peer type Phone <: { type Tie <: Multiple[Phone] & Single[Cloud] & Multiple[Edge] }(*@\label{line:phone-peer}@*)
\end{lstlisting}

Each phone is tied to exactly one cloud instance and to many other phones (\cref{line:phone-peer});
the cloud is tied to many phones and to one edge server (\cref{line:cloud-peer}).
%
The compiler rejects any access the architecture does not allow.

\paragraph{Crowding monitoring.}
Crowding is a property of the whole ensemble of phones, not of any two of them.
%
The ensemble's membership changes as participants arrive and leave,
and no device holds enough information to compute the estimate on its own.
%
In aggregate computing (\cref{subsec:aggregate-computing}) this becomes a field computation.
The program is written once for the ensemble, and every device runs it against its own neighbourhood.
%
The estimate uses three self-stabilising building blocks:
a gradient from a leader, a convergecast that sums the observations along the gradient, and a second convergecast that measures the radius of the region.
%
Dividing the sum by the area gives the density.

\paragraph{Point-of-interest recommendations.}
The recommendation flow is a protocol:
the phone presents an identifier, the cloud authorises it, the edge server produces a recommendation, and the phone displays it.
%
These exchanges must occur in this order.
%
Choreographic programming (\cref{subsec:choreographic-programming}) specifies that order once as a global description and projects it onto each participant,
so the endpoint code cannot deviate from it.

\subsection{The Cost of Combining Paradigms}
\label{subsec:multiparty-challenge}

Each paradigm fits its own part of \cityevt, but the three parts are not independent.
%
The recommendation protocol needs the identifiers produced by the multitier part,
and the organisers' dashboard needs the crowding level computed by the aggregate part.

Running three frameworks side by side does not preserve their guarantees across framework boundaries.
%
Each framework has its own abstractions, connectivity assumptions, and model of where computation happens,
and it checks its properties against that model.
%
No model governs interactions between the frameworks.
For example, a value leaves the choreography and enters the multitier code as an ordinary value,
with no check that the receiving peer may read it.
%
Code from another framework can also violate a paradigm's guarantees.
%
A choreography is deadlock-free because its projection fixes who sends and who receives.
%
Nesting one of its blocks inside a multitier scope that belongs to a different peer breaks that guarantee:
the send may never execute, and the receiver waits for a message no one will produce.

Two notions are therefore missing: where a value lives, and which operations a given piece of code may perform.

\section{Capabilities as a Unifying Abstraction}
\label{sec:capabilities-unification}

\camil models each paradigm as a capability~\cite{Farabegoli:2026:CaMiL}.
%
A capability is a value that carries the authority to perform a class of operations:
code that does not hold it cannot call them.
%
Authority comes from the structure of the program alone, and no ambient state can grant it.
%
The design draws on work discussed in \cref{subsec:direct-style-effect-systems}:
object capabilities as unforgeable references~\cite{miller2003capability},
capability passing as a language design~\cite{brachthaeuser2020effekt,10.1145/3428194},
and capture checking~\cite{DBLP:journals/toplas/BoruchGruszeckiOLLB23}.
%
Capture checking puts capability tracking in Scala's type system: what a value may use becomes part of its type.

\begin{table}[t]
    \centering
    \caption{The capabilities of CaMiL and the operations each one guards.}
    \label{tab:camil-capabilities}
    \begin{tabular}{ll}
        \toprule
        \textbf{Capability} & \textbf{Guarded operations} \\
        \midrule
        \texttt{Placement}   & \texttt{on[P]}, \texttt{take} \\
        \texttt{Multitier}   & \texttt{asLocal}, \texttt{asLocalAll} \\
        \texttt{Choreography} & \texttt{comm}, \texttt{multicast}, \texttt{broadcast} \\
        \texttt{Collective}  & \texttt{collective}, \texttt{rep}, \texttt{nbr}, \texttt{branch} \\
        \midrule
        \texttt{Scope[L]}, \texttt{PeerScope[P]} & scope tokens gating all of the above \\
        \bottomrule
    \end{tabular}
\end{table}

\camil has four capabilities: one for each of the three paradigms and \code{Placement}.
\Cref{tab:camil-capabilities} lists the operations guarded by each capability.
%
The prototype expresses them as Scala~3 \emph{contextual parameters}.
A method declares the capabilities it needs with the \code{using} keyword.
The type system resolves them from the surrounding lexical context.
%
An operation is then available only where its paradigm's capability is in scope,
so the three paradigms stay separate without relying on any coding convention.
%
All four capabilities share placement types:
one paradigm hands a value to another through them.

\subsection{Placement Types and the Placement Capability}
\label{subsec:placement-capability}

Placement types were introduced in multitier programming to specify where a value resides~\cite{DBLP:journals/pacmpl/WeisenburgerKS18}.
%
\camil uses them in all three paradigms, so location is written the same way everywhere.
%
The \code{Placement} capability provides two operations:
\code{on} places a value on a peer, and \code{take} reads a value placed on the local peer.

\begin{lstlisting}[
  style=scala3,
  escapeinside={(*@}{@*)},
  caption={The placement capability: placing a value on a peer and reading it back.},
  label={lst:placement-capability},
]
def helloOnSmartphone(using Placement): String on Phone = on[Phone]("Hello")(*@\label{line:placement-on}@*)
def reply(using Placement): String on Phone = on[Phone]:(*@\label{line:placement-take}@*)
  val hello: String = take(helloOnSmartphone)
  s"${(*@hello@*)}, Smartphone!"
\end{lstlisting}

In \cref{line:placement-on}, \code{on} places the string on the \code{Phone} peer,
so the result has type \code{String}\;\code{on}\;\code{Phone}.
%
In \cref{line:placement-take}, \code{take} reads the local value.
%
Both functions declare \code{Placement}, and without it neither operation is in scope.

The type system also checks that \code{take} runs on the peer where the value resides:
taking a value placed elsewhere does not compile, so no check is needed at run time.
%
\code{Placement} does not move values between peers.
It only names locations; the three paradigm capabilities move values between them.

\subsection{The Multitier Capability}
\label{subsec:multitier-capability}

The \code{Multitier} capability guards the two tier-crossing operations.
%
\code{asLocal} brings a value from a single remote peer to the local one;
\code{asLocalAll} aggregates the values of many remote peers into a map.
%
In multitier programming, communication may appear anywhere inside a placed scope that needs a remote value.
%
\Cref{lst:multitier-registration} is \cityevt's registration written with this capability.

\begin{lstlisting}[
  style=scala3,
  escapeinside={(*@}{@*)},
  caption={User registration with the multitier capability.},
  label={lst:multitier-registration},
]
def registration(userId: String)(using Placement, Multitier) =
  val registrationRequest: Form on Phone = on[Phone] { /* fill form */ }(*@\label{line:mt-phone}@*)
  val ids: Map[Remote[Phone], UUID] on Cloud = on[Cloud]:
    val request = asLocalAll(registrationRequest)(*@\label{line:mt-cloud}@*)
    request.view.mapValues(registerUser).toMap
  val id: UUID on Phone = on[Phone]:
    val localId = getId(localAddress)
    asLocal(ids).get(localId)(*@\label{line:mt-aslocal}@*)
\end{lstlisting}

The request is placed on the phone (\cref{line:mt-phone}).
%
The cloud collects the requests of every connected phone with \code{asLocalAll} and registers each one (\cref{line:mt-cloud}).
%
Each phone then reads back its own identifier with \code{asLocal} (\cref{line:mt-aslocal}).
%
The program contains no explicit send operation.
All communication uses these two operations,
which are checked against the ties of \cref{lst:scalaloci-architecture}.

\subsection{The Aggregate Capability}
\label{subsec:aggregate-capability}

The \code{Collective} capability provides \code{rep}, \code{nbr}, and \code{branch},
the field-calculus constructs~\cite{VIROLI2019100486} described in \cref{subsubsec:ac-execution-model}.
It also provides an entry point that opens a collective computation on a peer type and sets its round duration.

\begin{lstlisting}[
  style=scala3,
  escapeinside={(*@}{@*)},
  caption={Crowd density estimation with the collective capability.},
  label={lst:collective-crowding},
]
def crowding(using Placement, Collective): Signal[Double] on Phone =
  Collective[Phone](1.s):(*@\label{line:coll-scope}@*)
    // gradient from the leader of a 300-metre region
    val distanceToLeader = distanceTo(S(300))
    val people = /* people detected by local sensors */
    // convergecast the observations, and the extent of the region
    val count = C(distanceToLeader, people, _ + _)
    val radius = C(distanceToLeader, distanceToLeader, _ max _)
    count / (math.Pi * radius * radius)   // density = count / area
\end{lstlisting}

\Cref{line:coll-scope} opens a collective scope on the \code{Phone} peer type, with one-second rounds.
%
Inside, the estimate uses the same self-stabilising blocks available in ScaFi~\cite{casadei2022scafi}.
%
The result has type \code{Signal[Double]}\;\code{on}\;\code{Phone}:
a stream that emits a new density value at every round, placed on the phones that compute it.
%
Because the result carries a placement type, the other two paradigms can read it.

A collective computation also requires its peer type to be tied to itself with multiple cardinality.
%
Aggregate computing usually leaves this assumption implicit:
a field computation runs over a homogeneous ensemble of devices that can reach each other,
not over an arbitrary set of peers.

\subsection{The Choreography Capability}
\label{subsec:choreography-capability}

The \code{Choreography} capability provides \code{comm}, \code{multicast}, and \code{broadcast}.
%
\code{comm} moves a value between one sender and one receiver and requires a single tie in both directions;
\code{multicast} sends from one peer to many;
\code{broadcast} makes a value available to all peers.
%
\Cref{lst:choreography-recommendation} is the recommendation protocol.

\begin{lstlisting}[
  style=scala3,
  escapeinside={(*@}{@*)},
  caption={The recommendation protocol with the choreography capability.},
  label={lst:choreography-recommendation},
]
def retrievePreference(using Placement, Choreography) = Choreography:(*@\label{line:chor-scope}@*)
  val smartphoneToken = on[Phone] { /* get id */ }
  val cloudToken = comm[Phone, Cloud](smartphoneToken)(*@\label{line:chor-comm1}@*)
  val isRegistered = on[Cloud] { take(cloudToken) /* registered? */ }
  val shouldRecommend = comm[Cloud, Edge](isRegistered)
  val recommendation = on[Edge]:
    if take(shouldRecommend) then { /* provide recommendation */ }
  val recommOnSmartphone = comm[Edge, Phone](recommendation)
  on[Phone] { take(recommOnSmartphone) /* use recommendation */ }
\end{lstlisting}

The protocol reads top to bottom as one global description:
the identifier travels from phone to cloud (\cref{line:chor-comm1}),
the authorisation from cloud to edge,
and the recommendation back to the phone.
%
Each \code{comm} names both endpoints and is checked against the architecture.
%
Unlike the multitier code of \cref{lst:multitier-registration}, communication forms the structure of this program.
The placed computations are attached to that structure.

\subsection{Composing Paradigms}
\label{subsec:paradigm-composition}

Each capability opens a scope where only its own operations are available.
%
Operations never cross between scopes, but values do, because all the scopes share the same placement types.
%
Multitier code reads a remote placed value with \code{asLocal}.
Choreographic code moves one between peers with \code{comm}.
A collective computation produces one placed value per round, which the other two paradigms can read.
%
The architecture, not the capability that produced the value, decides who may read it.

\Cref{lst:composition-recommendation} combines choreographic and multitier code in \cityevt's recommendation workflow.

\begin{lstlisting}[
  style=scala3,
  escapeinside={(*@}{@*)},
  caption={CityEvt's recommendation workflow, combining choreographic and multitier code.},
  label={lst:composition-recommendation},
]
object RecommendationSystem:
  def recomm(smartphoneToken: Token on Phone)(using
    Placement, Choreography, Multitier(*@\label{line:comp-caps}@*)
  ) = Multitier:
    val authRequests = on[Edge] { asLocalAll(smartphoneToken) }
    val authPeers = Choreography:(*@\label{line:comp-chor}@*)
      val requests = comm[Edge, Cloud](authRequests)(*@\label{line:comp-comm}@*)
      val authGranted = on[Cloud] { /* grant access */ }
      comm[Cloud, Edge](authGranted)
    val recommendations = on[Edge]:
      createRecommendations(take(authPeers))(*@\label{line:comp-take}@*)
    on[Phone]:
      asLocal(recommendations).subscribe: recommendation => /* display */(*@\label{line:comp-aslocal}@*)
\end{lstlisting}

The function declares the three capabilities it needs (\cref{line:comp-caps}), which makes their operations available in its body.
%
A choreographic scope sits inside a multitier one (\cref{line:comp-chor}).
%
Two values cross the boundary in opposite directions.
The authentication requests gathered by the multitier code enter the choreography (\cref{line:comp-comm}).
The multitier code then reads the result of the choreography (\cref{line:comp-take}).
%
Both values are ordinary placed values and are type-checked as such.

\Cref{lst:composition-crowding} combines collective and multitier code.

\begin{lstlisting}[
  style=scala3,
  escapeinside={(*@}{@*)},
  caption={Reporting the collectively computed crowding level to the cloud.},
  label={lst:composition-crowding},
]
def crowdingLevelAndCloud(using Placement, Collective, Multitier) =
  val level: Signal[Double] on Phone = Collective[Phone](1.s):(*@\label{line:comp-collective}@*)
    /* compute crowding level */
  Multitier:
    on[Cloud]:
      asLocalAll(level).values.foreach: signal =>(*@\label{line:comp-collect-all}@*)
        signal.subscribe: level =>
          /* monitor crowding level on the cloud */
\end{lstlisting}

The collective scope produces a signal placed on the phones (\cref{line:comp-collective}),
each of which computes its own local estimate.
%
The multitier scope then gathers those signals at the cloud with \code{asLocalAll} and subscribes to each of them (\cref{line:comp-collect-all}).
%
The two scopes run on different execution models:
the collective computation proceeds in rounds, while the multitier code reads remote values whenever it needs them.
%
The placed signal connects the two:
the collective side decides when a new value appears, and the multitier side decides who reads it.

\subsection{Safety Properties}
\label{subsec:capability-safety}

The capability discipline rules out three errors that appear when the paradigms are combined by hand.

\paragraph{Safe boundaries.}
In \cref{lst:composition-recommendation,lst:composition-crowding} the paradigms meet only at placed values.
%
Inside a scope, the type system enforces the rules of that paradigm;
between scopes, the only channel is a placed value, whose type says where it lives.
%
The paradigms cannot interact directly outside these checked boundaries.

\paragraph{Leak prevention.}
A closure can capture a placed value.
If the closure is sent to another peer and invoked there,
it tries to read a value that does not exist on that peer.
%
\Cref{lst:leak-prevention} shows the pattern.

\begin{lstlisting}[
  style=scala3,
  escapeinside={(*@}{@*)},
  caption={A placed closure that captures a local value and is then sent to another peer.},
  label={lst:leak-prevention},
]
def choreograph(*@@*)y(using Placement, Choreography): Unit = Choreography:
  val valueOnSmartphone: String on Phone = on[Phone]:(*@\label{line:leak-value}@*)
    "Hello from Smartphone!"
  val funAtSmartphone: (() => String) on Phone = on[Phone]:(*@\label{line:leak-closure}@*)
    () => take(valueOnSmartphone)
  val msgAtCloud: (() => String) on Cloud = comm[Phone, Cloud](funAtSmartphone)(*@\label{line:leak-comm}@*)
  on[Cloud]:
    val f: () => String = take(msgAtCloud)
    f()(*@\label{line:leak-call}@*)
\end{lstlisting}

A string is placed on the phone (\cref{line:leak-value}) and captured by a function also placed there (\cref{line:leak-closure}).
%
The function is sent to the cloud (\cref{line:leak-comm}) and invoked there (\cref{line:leak-call}).
The captured string is not reachable from the cloud.
%
The closure's type tracks the peer-specific capability its \code{take} needs,
so \camil rejects the program at compile time rather than letting the \code{take} fail at run time.

\paragraph{Deadlock avoidance.}
An \code{on} block runs its body only on the peer its type names, and skips it everywhere else,
including any block nested inside it.
%
If two nested \code{on} blocks name different peers, the inner body never runs on any peer.
The system deadlocks when a participant waits for a value that the inner body should have produced.

\begin{lstlisting}[
  style=scala3,
  escapeinside={(*@}{@*)},
  caption={Nested placed blocks naming different peers, a pattern that causes deadlock.},
  label={lst:deadlock-avoidance},
]
def choreograph(*@@*)y(using Placement, Choreography): String = Choreography:
  val msgAtSmartphone: String on Phone = on[Phone]:(*@\label{line:dl-outer}@*)
    val onCloud: String on Cloud = on[Cloud]:(*@\label{line:dl-inner}@*)
      "Hello from Cloud!"
    val onSmartphone: String on Phone = comm[Cloud, Phone](onCloud)(*@\label{line:dl-comm}@*)
    take(onSmartphone)
\end{lstlisting}

The cloud never runs the inner block of \cref{line:dl-inner},
because it sits inside a block the cloud skips (\cref{line:dl-outer}).
%
The phone then waits on \cref{line:dl-comm} for a value that is never produced.
%
\camil also rejects this program at compile time through the mechanism described in \cref{subsec:scala3-machinery}.

\section{Realisation}
\label{sec:camil-realisation}

\subsection{Type-level Machinery in Scala 3}
\label{subsec:scala3-machinery}

\camil needs three features from its host language, and Scala~3 provides them.

Contextual abstractions declare that an operation requires a capability.
%
A context parameter, written \code{using}, makes the capability an argument of the operation;
a context function, written \code{?=>}, supplies it to a block.
%
In both cases the compiler resolves the argument at the call site from the lexical context.
%
\camil gives each paradigm a marker for its own scope and gates that paradigm's operations on the marker.

\begin{lstlisting}[
  style=scala3,
  escapeinside={(*@}{@*)},
  caption={The base capability and the per-paradigm scope capability.},
  label={lst:scope-capabilities},
]
trait Multiparty extends SharedCapability
trait Scope[L] extends SharedCapability
\end{lstlisting}

\code{Multiparty} is the base capability for the whole framework;
\code{Scope[L]} marks one paradigm's lexical scope.
%
Each paradigm's companion object then supplies an \code{apply} that opens its scope,
as \cref{lst:multitier-object} shows for the multitier case.

\begin{lstlisting}[
  style=scala3,
  escapeinside={(*@}{@*)},
  caption={The multitier capability: a scope marker, an operation gated on it, and the entry point that supplies it.},
  label={lst:multitier-object},
]
object Multitier:
  sealed trait MultitierScope extends Erased(*@\label{line:mt-scope}@*)

  def asLocal[L <: TiedSingleWith[R], R <: Peer: PeerTag, V, SC](
    using Multitier, PeerScope[L], Scope[SC], SC =:= MultitierScope)((*@\label{line:mt-evidence}@*)
    value: V on R): V

  def apply[A](using Multitier)(multitie(*@@*)r: Scope[MultitierScope] ?=> A): A(*@\label{line:mt-apply}@*)
  ... // other operations
\end{lstlisting}

\Cref{line:mt-scope} declares the scope marker.
%
It extends \code{Erased}, so it exists only for type-level tracking and has no representation at run time.
%
Each operation asks for a \code{Scope[SC]} together with evidence that \code{SC} is \code{MultitierScope} (\cref{line:mt-evidence}),
and this confines it to the innermost multitier scope.
%
\code{apply} is the only way to obtain that capability (\cref{line:mt-apply}).
It passes the capability to the block through a context function and itself requires \code{Multitier}.
%
\code{Collective} and \code{Choreography} have the same shape.

Capture checking keeps a capability from escaping its scope
by tracking captured capabilities in types~\cite{DBLP:journals/toplas/BoruchGruszeckiOLLB23}.
%
\camil uses it twice.
%
Because \code{Scope[L]} is a capability, the compiler does not let it leave the block passed to \code{apply}.
Returning it is a capture error,
so one paradigm's authority cannot be carried outside its scope.
%
The same mechanism rules out the leak of \cref{lst:leak-prevention}.
The peer-specific capability required by the closure appears in its type.
Sending the closure to a peer that cannot supply that capability therefore does not type-check.

Separation checking keeps a capability from being aliased.
%
It distinguishes read-only capabilities, which can be shared freely, from capabilities with side effects, which cannot~\cite{DBLP:journals/pacmpl/XuBO24}.
%
\camil declares \code{Placement} mutable, as \cref{lst:placement-mutable} shows.

\begin{lstlisting}[
  style=scala3,
  escapeinside={(*@}{@*)},
  caption={The placement capability is mutable, so it cannot be aliased.},
  label={lst:placement-mutable},
]
trait Placement extends Mutable:
  update def on[P <: Peer][V](body: PeerScope[P] ?=> V): V on P
\end{lstlisting}

\code{on} mutates the capability, since it produces a fresh identifier for every placed value.
%
Separation checking forbids aliasing a mutable capability.
A nested \code{on} would need a second reference to \code{Placement} while the enclosing one still holds the first.
%
Nesting therefore does not compile.
%
The deadlock of \cref{lst:deadlock-avoidance} is ruled out by that general rule, and not by a check written for it.

\subsection{Formal Account}
\label{subsec:camil-formalization}

A core calculus gives the four capabilities a single semantics and defines the meaning of programs that mix paradigms~\cite{Farabegoli:2026:CaMiL}.
%
As in \cref{subsec:pulverization-execution-model}, the full rules are given in the paper rather than repeated here.

The expression language extends the simply-typed lambda calculus with the placement expressions \code{on} and \code{take}.
It also provides a \code{handle} construct that supplies a capability and starts execution.
%
Function types record a capability set:
calling a function of type $\mptype[1] \rightarrow_{\capset} \mptype[2]$ incurs the capabilities in $\capset$,
and an empty set means the function is pure.
%
The set contains the four capabilities and the peer capabilities, which tie an operation to one peer type.
%
The peers are those declared in the architecture specification,
and a placement type $\ton{\mptype}{\pfam}$ pairs a type with the peer that holds its values.

Two judgements give the semantics.
%
Typing is annotated with the capabilities an expression needs,
\[
    \typing{\Gamma}{\mpexpr}{\mptype}{\capset},
\]
which reads: under the architecture $\arch$ and the environment $\Gamma$, the expression $\mpexpr$ has type $\mptype$ and needs the capabilities $\capset$.
%
Reduction is annotated with the peer instance that performs the step,
\[
    \redn{\peerconf;\vmstate}{\mpexpr}{\peerconf';\vmstate'}{\mpexpr'}{\peerinst[i]{\pfam}},
\]
where $\peerconf$ maps each peer type to its live instances and $\vmstate$ holds the state of the aggregate virtual machine.
%
The semantics models a distributed run as one global execution thread annotated with peer instances.
Multitier languages use the same approach~\cite{DBLP:journals/pacmpl/WeisenburgerKS18,Zakhour:2023:TypeSafe}.

Each paradigm adds its own constructs and rules to this core.
Multitier programming adds \code{asLocal} and \code{asLocalAll}.
Choreography adds \code{comm}, \code{multicast}, and \code{broadcast}.
Aggregate computing adds \code{collective}, \code{rep}, \code{nbr}, and \code{branch}.
%
The aggregate constructs read and update $\vmstate$.
%
The typing rule for \code{on} removes the placement capability from the context of the body and adds the capability of the peer being entered.
%
A peer's capability therefore exists only inside the block that enters that peer,
and cannot appear in what the block returns.
%
This is why the closure of \cref{lst:leak-prevention} does not type-check.

Soundness holds under two conditions.
%
The peer configuration must satisfy the architecture:
every peer type has at least one instance, every single tie is realised by exactly one link, and every multiple tie by at least one.
%
The virtual-machine state must be consistent with the collective computations in progress.

\begin{theorem}[Soundness]
    \label{thm:camil-soundness}
    Let $\peerconf \vDash \arch$ be an architecture-satisfying peer configuration with consistent virtual-machine states,
    and let $\mpexpr$ be an expression at a peer instance $\peerinst[i]{\pfam}$ with $\typing{\Gamma}{\mpexpr}{\mptype}{\capset}$.
    Then $\mpexpr$ is either a value or it reduces to an expression of the same type that requires no additional capabilities.
    The virtual-machine states remain consistent.
\end{theorem}

The proof follows the standard progress-and-preservation argument by induction on the typing and reduction derivations.
%
The theorem establishes two properties.
%
Every \code{comm}, \code{asLocal}, or \code{asLocalAll} step travels along a tie declared by the architecture.
The typing rule for each operation requires that tie as a premise.
A well-typed program therefore conforms to its architecture.
%
Because the capability sets of the three paradigms are disjoint,
no well-typed program runs one paradigm's operation under another's rules.
%
The paradigms can interact only through placed values,
and this is why the composition of \cref{subsec:paradigm-composition} is safe:
no pairwise compatibility check between paradigms is needed.

\section{Evaluation}
\label{sec:camil-evaluation}

\subsection{Expressivity and Safety}
\label{subsec:camil-expressivity}

The expressivity evaluation asks whether one capability-based framework can express programs supported by three dedicated languages.
%
The 46 use cases of \cref{tab:camil-use-cases} come from the multitier, choreographic, and aggregate literature and cover seven application classes;
all of them were reimplemented in \camil~\cite{Farabegoli:2026:CaMiL}.

\begin{table}[t]
    \centering
    \caption{Use cases reimplemented in CaMiL, with the paradigm each one exercises.}
    \label{tab:camil-use-cases}
    \scriptsize
    \begin{minipage}[t]{0.49\textwidth}
        \vspace{0pt}
        \begin{tabular}[t]{@{}llc@{}}
            \toprule
            \textbf{Class} & \textbf{Use case} & \textbf{Parad.} \\
            \midrule
            Communic. & Basic Chat & \emph{Multi.} \\
                      & Chat Application & \emph{Multi.} \\
                      & Chirper & \emph{Multi.} \\
                      & Chirper Inc & \emph{Multi.} \\
                      & Haste Chatbox & \emph{Multi.} \\
                      & Email Application & \emph{Chor.} \\
                      & Email System & \emph{Multi.} \\
                      & Vitals Streaming & \emph{Chor.} \\
                      & Ping-Pong & \emph{Chor.} \\
                      & Multicast & \emph{Chor.} \\
                      & Broadcasting & \emph{Aggr.} \\
                      & Channel & \emph{Aggr.} \\
            \midrule
            Coordin.  & Book Seller & \emph{Chor.} \\
                      & Buyer-Seller-Shipper & \emph{Chor.} \\
                      & Fan-in & \emph{Chor.} \\
                      & Fan-out & \emph{Chor.} \\
                      & GMW & \emph{Chor.} \\
                      & Main-Worker & \emph{Multi.} \\
                      & Token Ring & \emph{Multi.} \\
                      & Remote Function & \emph{Chor.} \\
                      & Hello Roles & \emph{Chor.} \\
            \midrule
            Security  & Diffie-Hellman & \emph{Chor.} \\
                      & Distributed Auth. & \emph{Chor.} \\
                      & DPrio & \emph{Chor.} \\
                      & Ticket Trader & \emph{Multi.} \\
            \bottomrule
        \end{tabular}
    \end{minipage}
    \hfill
    \begin{minipage}[t]{0.49\textwidth}
        \vspace{0pt}
        \begin{tabular}[t]{@{}llc@{}}
            \toprule
            \textbf{Class} & \textbf{Use case} & \textbf{Parad.} \\
            \midrule
            Games       & Card Game & \emph{Chor.} \\
                        & Lottery & \emph{Chor.} \\
                        & Tic-Tac-Toe & \emph{Chor.} \\
                        & Guess a Number & \emph{Multi.} \\
            \midrule
            Data/Serv.  & Key-value Store & \emph{Chor.} \\
                        & Shopping List & \emph{Multi.} \\
                        & Graffiti RPC & \emph{Multi.} \\
                        & Graffiti WebSockets & \emph{Multi.} \\
                        & Tweet Processing & \emph{Multi.} \\
                        & Consume Items & \emph{Chor.} \\
            \midrule
            Algo./Comp. & Karatsuba & \emph{Chor.} \\
                        & Merge Sort & \emph{Chor.} \\
                        & Quicksort & \emph{Chor.} \\
                        & Collection & \emph{Aggr.} \\
                        & Collection IDs & \emph{Aggr.} \\
                        & Count Neighbours & \emph{Aggr.} \\
                        & Gradient Explicit Fields & \emph{Aggr.} \\
                        & Max ID & \emph{Aggr.} \\
                        & Neighbor Counter & \emph{Aggr.} \\
            \midrule
            Spatial     & Neighbor Distances & \emph{Aggr.} \\
                        & Neighbor Dist. Obstacle & \emph{Aggr.} \\
            \bottomrule
        \end{tabular}
    \end{minipage}
\end{table}

The reimplementations total 3874 lines of \camil code:
2012 for the choreographic use cases, 1425 for the multitier ones, and 437 for the aggregate ones.
%
For this corpus, representing a paradigm as a capability therefore did not reduce expressivity relative to a dedicated language.

The same use cases test the safety properties of \cref{subsec:capability-safety}.
%
For each paradigm, \camil rejects programs that are ill-formed under the rules of the corresponding dedicated language.
It rejects communication along undeclared ties and attempts to take a placed value on the wrong peer.
It also rejects collective computations over peers that are not tied to each other.
%
Each dedicated language already performs its own checks;
what \camil adds is one type system that performs all of them.

\subsection{Inter-paradigm Integration}
\label{subsec:inter-paradigm}

The evaluation also tests whether the paradigms can work together in one program.
%
\Cref{tab:camil-integration} summarises three applications that mix them.

\begin{table}[t]
    \centering
    \caption{Applications combining several paradigms and the contribution of each paradigm.}
    \label{tab:camil-integration}
    \scriptsize
    \setlength{\tabcolsep}{4pt}
    \begin{tabularx}{\textwidth}{@{}l>{\raggedright\arraybackslash}X>{\raggedright\arraybackslash}X>{\raggedright\arraybackslash}X@{}}
        \toprule
        \textbf{Application} & \textbf{Multitier} & \textbf{Choreographic} & \textbf{Collective} \\
        \midrule
        CityEvt & User registration and recommendation delivery across phones, edge, and cloud. & Recommendation protocol between edge nodes and the cloud. & Crowding-level monitoring from nearby phones. \\
        \addlinespace
        SmrtMob & Vehicle registration, credential distribution, and edge-side traffic-light coordination. & Route recommendation protocol across vehicles, edge nodes, and the cloud. & Traffic sensing from nearby vehicles to derive local traffic probes. \\
        \addlinespace
        Nebula & Dashboard reporting of learning progress and network status. & Round coordination for training, aggregation, and lifecycle management. & Distance-based neighbour selection for topology-aware participation. \\
        \bottomrule
    \end{tabularx}
\end{table}

\cityevt has three peer families and three subsystems.
The paradigms exchange values at two points.
%
\smrtmob was written from scratch to test this combination.
Vehicles sense roadway conditions while moving.
Edge nodes summarise local congestion and coordinate traffic lights,
while the cloud computes route recommendations from a wider snapshot.
%
Its four phases use three capabilities:
registration and traffic-light coordination are multitier,
traffic sensing is collective,
and route recommendation is choreographic.
%
The implementation is about 500 lines long and has five boundaries where one paradigm reads a value produced by another.

Nebula tests integration in a system that was not designed around any of these paradigms.
%
It is an open-source platform for decentralised federated learning with a public codebase of about 132\,000 lines~\cite{beltran2023decentralized,MartinezBeltran:fedstellar:2024,Nebula}.
The platform consists of a dashboard, an orchestrating controller, and a learning engine that runs on distributed nodes.
%
Three of its parts were reimplemented in \camil:
round coordination as a choreographic protocol between the coordinator and the participants,
the distance-based neighbour policy as an aggregate program,
and dashboard reporting as multitier code.
%
The aggregate layer computes topology information:
the choreographic protocol uses it to decide who participates in a round, and the dashboard uses it to report status.
%
The choreographic layer, in turn, produces the global model and the round summaries that the dashboard displays.

The reimplementations remove explicit message-passing code from each part.
%
The original implements the round lifecycle across 22 message callbacks.
In \camil, one choreographic protocol replaces them without callbacks or inversion of control.
%
The original neighbour policy interacts with the communication infrastructure at 57 points.
In \camil, it becomes a single aggregate program that leaves neighbour discovery and message exchange to the collective runtime.
%
The networking layer defines 60 message and action types for moving data between hosts.
The reimplementation uses seven communication terms.
%
These numbers count code, not run time.

\section{Monadic Communication Primitives}
\label{sec:scalatropy}

\camil unifies the three paradigms through capabilities.
A paradigm's operations are available where its capability is in scope.
Placement types carry values across paradigm boundaries.
%
ScalaTropy covers two of them, choreographic and multitier programming, and uses a different mechanism~\cite{Farabegoli:2026:ScalaTropy}:
the monadic sequencing described in \cref{subsec:monadic-effects}.
%
A coordination program is a pipeline of communication steps composed with bind.
The selected interpreter executes those steps as either real network exchanges or a local simulation.
%
ScalaTropy also addresses a limitation shared by both paradigms.
They provide point-to-point communication but no dedicated operation for sending a different message to each recipient.
This pattern must instead be encoded manually, usually through broadcasting and filtering.

\subsection{Communication Patterns in Multiparty Languages}
\label{subsec:communication-patterns}

An architecture follows the multitier tradition.
It consists of peer types and ties between them, each with a cardinality.
%
Writing $\pfam$ for a peer type, an architecture is generated by
\[
    \arch ::= \{\pfam\} \;\mid\; \arch \cup \arch \;\mid\; \tieSingle{\pfam[1]}{\pfam[2]} \;\mid\; \tieMult{\pfam[3]}{\pfam[4]},
\]
where $\tieSingle{\pfam[1]}{\pfam[2]}$ is a tie to exactly one instance and $\tieMult{\pfam[3]}{\pfam[4]}$ a tie to many.
%
At run time a configuration $\peerconf$ records the live instances $\peerinst[n]{\pfam}$ of each peer type and the links between them.

\begin{definition}[Architecture satisfaction]
    \label{def:arch-satisfaction}
    A configuration $\peerconf$ satisfies an architecture $\arch$, written $\peerconf \vDash \arch$, when
    every peer type of $\arch$ has at least one instance in $\peerconf$,
    every tie $\tieSingle{\pfam[1]}{\pfam[2]}$ is realised by exactly one link from each instance of $\pfam[1]$,
    and every tie $\tieMult{\pfam[3]}{\pfam[4]}$ by at least one link from each instance of $\pfam[3]$.
\end{definition}

In a satisfying configuration, four patterns exhaust the ways a message can be addressed.

\begin{definition}[Point-to-point communication]
    \label{def:point-to-point}
    One sender transmits a message $m$ to one receiver: $\peerinst{\pfam[s]} \xrightarrow{m} \peerinst{\pfam[r]}$.
\end{definition}

\begin{definition}[Isotropic communication]
    \label{def:isotropic}
    One sender transmits the \emph{same} message $m$ to every instance $\peerinst[1]{\pfam[r]},\dots,\peerinst[n]{\pfam[r]}$ of the receiving peer type.
\end{definition}

\begin{definition}[Anisotropic communication]
    \label{def:anisotropic}
    One sender transmits a \emph{different} message $m_i$ to each instance $\peerinst[i]{\pfam[r]}$ of the receiving peer type.
\end{definition}

\begin{definition}[Co-anisotropic communication]
    \label{def:co-anisotropic}
    Each instance $\peerinst[i]{\pfam[s]}$ of the sending peer type transmits its own message $m_i$ to a single receiver.
\end{definition}

\begin{figure}[t]
    \centering
    \begin{subfigure}[b]{0.23\textwidth}
        \centering
        \includegraphics[width=\textwidth]{point-to-point.pdf}
        \caption{Point-to-point}
        \label{fig:comm-point-to-point}
    \end{subfigure}
    \hfill
    \begin{subfigure}[b]{0.23\textwidth}
        \centering
        \includegraphics[width=\textwidth]{isotropic-comm.pdf}
        \caption{Isotropic}
        \label{fig:comm-isotropic}
    \end{subfigure}
    \hfill
    \begin{subfigure}[b]{0.23\textwidth}
        \centering
        \def\svgwidth{\textwidth}
        \input{figures/anisotropic-comm.pdf_tex}
        \caption{Anisotropic}
        \label{fig:comm-anisotropic}
    \end{subfigure}
    \hfill
    \begin{subfigure}[b]{0.23\textwidth}
        \centering
        \def\svgwidth{\textwidth}
        \input{figures/coanisotropic-comm.pdf_tex}
        \caption{Co-anisotropic}
        \label{fig:comm-co-anisotropic}
    \end{subfigure}
    \caption[Communication patterns in multiparty languages]{The four communication patterns. Isotropic communication sends one message to many receivers; anisotropic communication sends a different message to each; co-anisotropic communication collects a different message from each sender. Adapted from~\cite{Farabegoli:2026:ScalaTropy}.}
    \label{fig:comm-patterns}
\end{figure}

\Cref{fig:comm-patterns} shows the four patterns.
%
In isotropic communication, every recipient receives the same message.
In anisotropic communication, each recipient receives a distinct message.
%
A leader sending one configuration update to all replicas uses isotropic communication.
A coordinator assigning a distinct task to each worker uses anisotropic communication.
%
Sensors reporting their own readings to an aggregator are the co-anisotropic case, the fan-in dual of the anisotropic one.

\begin{table}[t]
    \centering
    \caption{Communication operations of three multiparty languages.}
    \label{tab:comm-operations}
    \footnotesize
    \begin{tabularx}{\textwidth}{@{}Xlll@{}}
        \toprule
        \textbf{Pattern} & \textbf{HasChor} & \textbf{ScalaLoci} & \textbf{ScalaTropy} \\
        \midrule
        Point-to-point & \texttt{\char`\~>} & \texttt{asLocal} & \texttt{comm} \\
        Isotropic      & --- & \texttt{asLocalFromAllSeq} & \texttt{isotropicComm} \\
        Anisotropic    & --- & \texttt{V per P on R} & \texttt{anisotropicComm} \\
        Co-anisotropic & --- & \texttt{V per R on P} & \texttt{coAnisotropicComm} \\
        \bottomrule
    \end{tabularx}
\end{table}

\Cref{tab:comm-operations} compares how three languages express these patterns.
%
HasChor, a monadic choreographic language in Haskell, provides only the point-to-point case~\cite{10.1145/3607849}.
%
ScalaLoci covers all four, but expresses the three non-point-to-point patterns at the type level rather than with dedicated operations.
%
An anisotropic assignment is a value of type \code{Signal[Task] per Worker on Master}, built by a function from a remote reference to a signal.
%
The reader has to recover the pattern from the type.
%
ScalaTropy gives each pattern its own primitive.

\subsection{The ScalaTropy DSL}
\label{subsec:scalatropy-dsl}

ScalaTropy distinguishes two forms of architecture.
%
Choreographic languages typically assume a \emph{weak} architecture, in which any participant may address any other;
multitier languages assume a \emph{strong} one, in which the ties fix who may talk to whom.
%
ScalaTropy uses the strong form and encodes it with Scala~3 structural types.

\begin{lstlisting}[
  style=scala3,
  escapeinside={(*@}{@*)},
  caption={A three-peer architecture as a type-level specification.},
  label={lst:scalatropy-architecture},
]
// Clients can communicate with a single Server
type Client <: { type Tie <: Single[Server] }
// Servers can communicate with a single Database and multiple Clients
type Server <: { type Tie <: Single[Database] & Multiple[Client] }
// The Database can communicate with one Server
type Database <: { type Tie <: Single[Server] }
\end{lstlisting}

The cardinalities of \cref{def:arch-satisfaction} become an enumeration indexed by the peer type,
and each cardinality gets a named refinement.

\begin{lstlisting}[
  style=scala3,
  escapeinside={(*@}{@*)},
  caption={Tie cardinalities and the refinements built from them.},
  label={lst:scalatropy-ties},
]
enum Quantifier[-P <: Peer]:
  case Single()
  case Multiple()

type TiedSingle[P <: Peer] = { type Tie <: Single[P] }
type TiedMultiple[P <: Peer] = { type Tie <: Multiple[P] }
\end{lstlisting}

As in \camil, placement types come from the multitier tradition.
Choreographies use located values, whereas multitier programs use placed values.
ScalaTropy represents both with ScalaLoci's infix \code{V}\;\code{on}\;\code{P}.
%
A value of type \code{V}\;\code{on}\;\code{P} is a value of type \code{V} residing at the instances of peer type \code{P}.

The DSL defines every operation in a single trait, parameterised by the monad \code{F} that sequences the computation.
%
\Cref{lst:scalatropy-multiparty} collects these primitives.

\begin{lstlisting}[
  style=scala3,
  escapeinside={(*@}{@*)},
  caption={The MultiParty trait: placed computation and the four communication primitives.},
  label={lst:scalatropy-multiparty},
]
trait MultiParty[F[_]]:
  type Remote[P <: Peer]
  type Anisotropic[P <: Peer, V]

  // Placed computation
  def on[L <: Peer, V](body: Label[L] ?=> F[V]): F[V on L](*@\label{line:st-on}@*)
  def take[L <: Peer](using Label[L])[V](placed: V on L): F[V]
  def takeAll[L <: Peer, R <: Peer](using Label[L])[V](
    placed: Anisotropic[R, V] on L): F[Map[Remote[R], V]]
  def reachablePeers[R <: Peer]: F[NonEmptyList[Remote[R]]](*@\label{line:st-reachable}@*)

  // Communication
  def comm[S <: TiedSingle[R], R <: TiedSingle[S], V](value: V on S): F[V on R](*@\label{line:st-comm}@*)
  def isotropicComm[S <: TiedMultiple[R], R <: TiedSingle[S], V](
    value: V on S): F[V on R](*@\label{line:st-isotropic}@*)
  def anisotropicMessage[S <: TiedMultiple[R], R <: TiedSingle[S], V](using Label[S])(
    value: Map[Remote[R], V], default: V): F[Anisotropic[R, V]](*@\label{line:st-aniso-msg}@*)
  def anisotropicComm[S <: TiedMultiple[R], R <: TiedSingle[S], V](
    value: Anisotropic[R, V] on S): F[V on R](*@\label{line:st-aniso}@*)
  def coAnisotropicComm[S <: TiedSingle[R], R <: TiedMultiple[S], V](
    value: V on S): F[Anisotropic[S, V] on R](*@\label{line:st-coaniso}@*)
\end{lstlisting}

The type bounds enforce the architecture.
%
\code{comm} (\cref{line:st-comm}) requires single-cardinality ties between the sender and receiver in both directions.
A point-to-point exchange between peers not paired by the architecture is therefore ill-typed.
%
\code{isotropicComm} (\cref{line:st-isotropic}) requires a multiple tie on the sender's side and a single tie on the receiver's.
These constraints enforce one-to-many communication.
%
The anisotropic case needs two steps.
\code{anisotropicMessage} (\cref{line:st-aniso-msg}) builds a per-recipient payload from a map,
using a default for recipients omitted from the map.
\code{anisotropicComm} (\cref{line:st-aniso}) then delivers each entry to its intended receiver.
%
\code{coAnisotropicComm} (\cref{line:st-coaniso}) is the dual, collecting a distinct value from each sender into one \code{Anisotropic} value at the receiver.
%
\code{reachablePeers} (\cref{line:st-reachable}) lets the program build that map without knowing the participants in advance.

\Cref{lst:scalatropy-isotropic} shows an isotropic exchange.

\begin{lstlisting}[
  style=scala3,
  escapeinside={(*@}{@*)},
  caption={An isotropic exchange that sends one update to every client.},
  label={lst:scalatropy-isotropic},
]
def isotropicExample[F[_]: Monad](using MultiParty[F]): F[Unit] = for
  updateFromServer <- on[Server] { "Update available." }
  updateAtClients <- isotropicComm[Server, Client](updateFromServer)(*@\label{line:st-iso-send}@*)
  _ <- on[Client]:
    take(updateAtClients).map(update => println(s"Client received: $update"))(*@\label{line:st-iso-recv}@*)
yield ()
\end{lstlisting}

The program is a \code{for} comprehension: a chain of binds in the monad \code{F}.
%
The single \code{isotropicComm} of \cref{line:st-iso-send} represents one transmission per client.
The receiving operation of \cref{line:st-iso-recv} is written once for all clients.
%
Nothing in the program names a participant instance:
the placement types name peer \emph{types}, and the instances are supplied at run time.

\subsection{Monadic Realisation}
\label{subsec:scalatropy-monad}

\Cref{lst:scalatropy-multiparty} leaves the monad \code{F} abstract.
This abstraction yields a tagless-final encoding~\cite{DBLP:journals/jfp/CaretteKS09}.
%
A program is written against the abstract interface, and an interpreter supplies the meaning of each operation.
%
Cloud Haskell's process monad uses the same separation for the same purpose (\cref{subsec:monadic-effects}).
The production interpreter runs the pipeline over a real network.
A test interpreter runs the same program as a pure in-memory computation,
so the two modes do not require separate code paths.
%
The alternative -- a free monad -- would require the semantics to be respecified for each interpreter;
tagless final avoids that at the cost of committing to a fixed set of operations.

A placed value is an algebraic data type with two constructors:
\code{Local}, carrying a deterministic identifier and the value itself,
and \code{Remote}, carrying only the identifier.
%
Executing \code{on[P]} compares \code{P} with the peer type of the host running the program:
on a match, the body runs and its result is wrapped as \code{Local};
otherwise the body is skipped and a \code{Remote} is returned.
%
When a communication primitive holds a \code{Local}, it sends the value.
When it holds a \code{Remote}, it retrieves the value from the network.
%
This design performs endpoint projection at run time rather than at compile time.
%
Census-polymorphic choreographic programming uses the same strategy to avoid fixing the participants at compile time~\cite{DBLP:journals/pacmpl/BatesKJSKN25}.

The interpreter rests on a small network interface:
send a value to an address, receive one from an address, report the live instances of a peer type, and name the local host.
%
The prototype implements this interface over MQTT, but the language does not depend on that choice;
HTTP or gRPC could be used instead.
%
The one requirement is that the live-instances query returns a non-empty list,
so that the anisotropic and co-anisotropic primitives can operate as topology membership changes.

\subsection{Case Studies and Communication Cost}
\label{subsec:scalatropy-evaluation}

\Cref{lst:scalatropy-master-worker} reimplements a standard example from the multitier literature in ScalaTropy~\cite{DBLP:journals/pacmpl/WeisenburgerKS18}.

\begin{lstlisting}[
  style=scala3,
  escapeinside={(*@}{@*)},
  caption={A master-worker application in ScalaTropy.},
  label={lst:scalatropy-master-worker},
]
object MasterWorker:
  type Master <: { type Tie <: Multiple[Worker] }
  type Worker <: { type Tie <: Single[Master] }

  case class Task(x: Int) { def compute: Int = x * x }

  def masterWorkerProgram[F[_]: MonadThrow](using MultiParty[F]): F[Unit] = for
    tasks <- on[Master]:
      for
        peers <- reachablePeers[Worker](*@\label{line:mw-peers}@*)
        allocation = peers.map(_ -> Task(Random.nextInt(100))).toList.toMap
        message <- anisotropicMessage[Master, Worker](allocation, Task(0))(*@\label{line:mw-message}@*)
      yield message
    taskOnWorker <- anisotropicComm[Master, Worker](tasks)(*@\label{line:mw-scatter}@*)
    partialResult <- on[Worker]:
      for
        t <- take(taskOnWorker)
        res <- t.compute.pure[F]
      yield res
    allResults <- coAnisotropicComm[Worker, Master](partialResult)(*@\label{line:mw-gather}@*)
    result <- on[Master] { takeAll(allResults).map(_.values.sum) }(*@\label{line:mw-reduce}@*)
  yield ()
\end{lstlisting}

The master discovers the currently reachable workers (\cref{line:mw-peers}) and builds one task per worker (\cref{line:mw-message}).
%
The scatter uses one \code{anisotropicComm} (\cref{line:mw-scatter}), and the gather uses one \code{coAnisotropicComm} (\cref{line:mw-gather}).
%
The gather also provides synchronisation.
The co-anisotropic value is complete only after every worker has contributed.
The master therefore cannot reduce it earlier (\cref{line:mw-reduce}), and the program needs no separate barrier.
%
A second case study ports a primary-backup replicated key-value store from HasChor.
It adds a third peer type and a tail-recursive loop over requests.
%
The ScalaTropy implementation works for any number of backups, while the choreographic implementation fixes the replica count statically~\cite{Farabegoli:2026:ScalaTropy}.

\begin{figure}[t]
    \centering
    \includegraphics[width=0.85\textwidth]{communication-overhead.pdf}
    \caption[Communication overhead of broadcast versus selective communication]{Bytes transmitted by the master for a matrix-vector product with increasing numbers of workers, using broadcast and selective communication. Data from~\cite{Farabegoli:2026:ScalaTropy}.}
    \label{fig:communication-overhead}
\end{figure}

Differentiated primitives affect both program structure and communication cost.
\Cref{fig:communication-overhead} reports their network cost.
%
The experiment computes a matrix-vector product over a fixed $50\times50$ matrix.
The master partitions the rows and distributes a distinct block to each worker.
%
The broadcast variant grows linearly with the number of workers,
because every worker receives the whole matrix and discards the other workers' blocks.
The selective variant remains nearly flat,
growing only with the per-worker metadata and the vector required by each worker.

Both variants are ScalaTropy implementations of the same application,
differing only in the primitive the scatter uses:
\code{isotropicComm} in one, \code{anisotropicComm} in the other.
%
The figure therefore measures the cost of encoding a selective pattern as broadcast-and-filter.
It does not compare ScalaTropy with HasChor or ScalaLoci, and it measures bytes rather than execution time.
%
The experiment does not measure two further properties.
Sending each recipient only the required data applies the principle of least privilege to communication.
Naming the pattern in the operation also makes the intent explicit,
rather than requiring the reader to infer it from the type.

\section{Discussion}
\label{sec:language-coordination-discussion}

The two designs address the same problem with different mechanisms.
%
In \camil, the basic notion is authority.
The capabilities held by a piece of code determine its permitted operations.
%
In ScalaTropy, the basic notion is sequencing.
A coordination program is a monadic pipeline whose interface determines the available operations.

They agree on the interface between paradigms.
Both use placement types from multitier programming and encode the architecture as ties with cardinalities.
%
Multitier programming supplies the vocabulary for the other two,
because it is the only one of the three paradigms of \cref{sec:foundational-paradigms} that already makes location explicit.

\camil covers all three paradigms and proves the combination sound.
%
ScalaTropy covers two paradigms and leaves the metatheory for future work.
However, it provides differentiated communication primitives that neither \camil nor the paradigms it unifies offer.
%
The designs are also complementary.
Nothing in \camil's capability design requires point-to-point communication to be the only choreographic primitive.
The four patterns of \cref{subsec:communication-patterns} could be added to its \code{Choreography} capability.

Both designs reject some programs accepted by the paradigms they unify.
%
\camil rejects nested placed scopes, which existing multitier languages allow.
%
This rule prevents deadlocks caused by nested scopes but also rejects safe nestings.
%
ScalaTropy's strong architecture imposes a related restriction.
A choreography that assumes full connectivity requires the developer to define an architecture first.

The pulverisation model of \cref{chap:pulverization} requires a macro-program represented as a graph of independently deployable components.
Each component has declared inputs and outputs.
%
The three paradigms produce such components, but in incompatible terms,
which is why a system built from more than one of them cannot be pulverised as a whole.
%
One notion of placement for all three removes that obstacle.
%
A placement type says where a component's value resides, and the architecture fixes which placements may exchange values.
%
The deployment model uses this information to determine where a component may run (\cref{sec:pulverization-abstractions}).
%
Neither design chooses a deployment.
The architecture only constrains the set of admissible deployments.
Runtime policies must choose among them and revise that choice as the infrastructure changes (\cref{chap:deployment-and-reconfiguration}).

Both designs still require the developer to write a program:
a capability-annotated Scala program, a monadic pipeline, a field program, or a choreography.
%
The remaining question is whether a macro-program can be generated from a natural-language specification (\cref{chap:language-based-macroprogramming}).
